\documentclass[aspectratio=169,14pt]{beamer}
\usepackage[utf8]{inputenc}\usepackage[T1]{fontenc}\usepackage[spanish]{babel}\usepackage{graphicx}
\usetheme{default}\definecolor{azul}{RGB}{35,91,133}\definecolor{rojo}{RGB}{162,39,50}
\setbeamercolor{structure}{fg=azul}\setbeamercolor{frametitle}{fg=azul}\setbeamercolor{alerted text}{fg=rojo}
\setbeamertemplate{navigation symbols}{}\setbeamertemplate{footline}{\hfill\scriptsize Aquiles Ramos Yapo | 227368 | R03\quad\insertframenumber/\inserttotalframenumber\quad\vspace{3mm}}
\setbeamertemplate{itemize items}[circle]\setbeamersize{text margin left=9mm,text margin right=9mm}
\newcommand{\dato}[1]{\csname dato#1\endcsname}\input{generado/valores}
\title{Minicentral hidroeléctrica en Limbani}\subtitle{Proyecto académico de diseño\\Sistema aislado para Limbani y Huancasayani}
\author{Estudiante: Aquiles Taylor Ramos Yapo\\Código: 227368 -- Grupo A}
\institute{Universidad Nacional del Altiplano\\Centrales Hidroeléctricas\\Docente: Mg. Ing. Gregorio Meza Marocho}
\date{14 de octubre de 2026}
\begin{document}
\begin{frame}\titlepage\end{frame}
\begin{frame}{Qué se diseña y qué se verificó}
\begin{columns}[T]\begin{column}{.50\textwidth}
\begin{itemize}\item Replanteamiento de una central referida como 200 kW, 50 m y Banki.
\item Inicialmente \textbf{aislada}; sin conexión al SEIN.
\item Se comparan 150, 200, 400 y 600 kW.
\end{itemize}
\end{column}\begin{column}{.46\textwidth}\includegraphics[width=\linewidth,height=4.1cm,keepaspectratio]{figuras/ubicacion-sandia.jpg}\par\scriptsize Sandia y Limbani. Imagen aportada en antecedentes; autor y fecha no identificados.\end{column}\end{columns}
\vspace{2mm}\small\alert{Proyecto estudiantil: cálculos de gabinete y comprobaciones de campo pendientes.}
\end{frame}
\begin{frame}{Ubicación: la central todavía necesita levantamiento}
\begin{columns}\begin{column}{.55\textwidth}\includegraphics[width=\linewidth,height=5.1cm,keepaspectratio]{figuras/emplazamiento-r03.png}\end{column}
\begin{column}{.42\textwidth}\small
Dos referencias de casa de máquinas difieren 152 m.\par\medskip
WGS84 / UTM 19S. La captación es una propuesta cartográfica.\par\medskip
\alert{No se sustituyen toma ni casa por las coordenadas del pueblo.}\par\medskip
Salto bruto de estudio: \dato{Hb} m, pendiente de nivelación.
\end{column}\end{columns}\end{frame}
\begin{frame}{Caudales: modelos públicos, no aforos locales}
\begin{columns}\begin{column}{.53\textwidth}\includegraphics[width=\linewidth,height=5.1cm,keepaspectratio]{figuras/caudales-r03.png}\end{column}
\begin{column}{.44\textwidth}\small
PISCO SENAMHI: 1981--2025.\par
16\,436 días y 540 meses.\par\medskip
Transposición por área geodésica:\par
\(\alpha=A_c/A_r=\dato{Alpha}\).\par\medskip
La estación Limbani publica niveles; no se recuperó curva de gasto.\par\medskip
\alert{Falta calibración en captación.}
\end{column}\end{columns}\end{frame}
\begin{frame}{El estiaje limita más que la potencia de placa}
\begin{columns}\begin{column}{.57\textwidth}\includegraphics[width=\linewidth,height=4.8cm,keepaspectratio]{figuras/duracion-r03.png}\end{column}
\begin{column}{.40\textwidth}\small
Q95 natural diario: \dato{Qlow} m\(^3\)/s.\par\medskip
Reservas de agua: supuestos, no autorizaciones.\par\medskip
Potencia diaria P95 en usuarios: \textbf{\dato{Pfirmstudy} kW}.\par
Mínimo del modelo: \textbf{0 kW}.\par\medskip
\alert{No se garantizan 150--200 kW continuos.}
\end{column}\end{columns}\end{frame}
\begin{frame}{Demanda: cuatro centros poblados, sin sumar el distrito}
\begin{itemize}
\item INEI 2017: Limbani y Huancasayani Centro, Bellavista y Lince.
\item \dato{PopBase} habitantes; \dato{VivsBase} viviendas ocupadas censales.
\item Proyección de estudio: \dato{Users} conexiones potenciales en 2026.
\item RD 108-2023: 600/400 W por lote y simultaneidad doméstica 0.80.
\end{itemize}
\medskip\large Pico de diseño: \textbf{\dato{Demand} kW}.\par
\small\alert{Vivienda ocupada no equivale a contrato activo. El padrón puede cambiar el dimensionamiento.}
\end{frame}
\begin{frame}{Alternativas y balance: cubrir pico exige respaldo}
\begin{columns}\begin{column}{.58\textwidth}\includegraphics[width=\linewidth,height=5cm,keepaspectratio]{figuras/oferta-demanda-r03.png}\end{column}
\begin{column}{.39\textwidth}\small
Caso conservador: 3 por 200 kW hidro.\par
Nominal en usuarios: \dato{Pusers} kW.\par\medskip
Energía hidro servida:\par\textbf{\dato{EHidro} MWh/año}.\par
Déficit: \dato{Edef} MWh/año.\par\medskip
Respaldo modular y gestión de carga.\par
\alert{600 kW no demuestra inversión óptima.}
\end{column}\end{columns}\end{frame}
\begin{frame}{Hidráulica: dimensiones coherentes con el cálculo}
\includegraphics[width=\textwidth,height=3.3cm,keepaspectratio]{figuras/obras-r03.png}\par
\small Salto bruto \dato{Hb} m; neto nominal \dato{Hn} m.\par
Canal 150 m; tubería 350 m, Di 1000 mm, espesor nominal 16 mm.\par
Caudal nominal \dato{Qnom} m\(^3\)/s; máximo de diseño 1.90.\par
\medskip\alert{Cotas de gabinete y trazado propuesto: falta comprobar la cámara en ladera.}
\end{frame}
\begin{frame}{Turbina, regulación y golpe de ariete}
\begin{columns}\begin{column}{.51\textwidth}\includegraphics[width=\linewidth,height=4.9cm,keepaspectratio]{figuras/transitorio-r03.png}\end{column}
\begin{column}{.46\textwidth}\small
Banki modular comparada con Francis y Pelton/Turgo.\par\medskip
Síncrono: 250 kVA, 380 V, 60 Hz, 900 rpm por grupo.\par\medskip
ELC rápido + gobierno lento + AVR coordinado.\par\medskip
Cierre de 8 s: \dato{PNslow} bar. Cierre brusco: cavitación posible.\par
\alert{PN40 es candidato, sin certificación del transitorio real.}
\end{column}\end{columns}\end{frame}
\begin{frame}{Electricidad: proteger una red aislada débil}
\includegraphics[width=\textwidth,height=3.9cm,keepaspectratio]{figuras/unifilar-r03.png}\par
\small T1 800 kVA, 10/0.38 kV; seis SED de 730 kVA en total.\par
Falla inicial máxima supuesta: \dato{IKmax} kA. La falla sostenida decrece.\par
\alert{Se necesita protección sensible a tensión y diferencial; no basta 50/51.}
\end{frame}
\begin{frame}{Seguridad y permisos: verificación explícita}
\begin{itemize}
\item CNE Suministro y Utilización con ámbitos distintos.
\item RD rural 2023, E.030 actualizada en 2026; matriz de 40 requisitos.
\item Tierra: resistividad, corriente y contacto/paso por verificar.
\item ANA: disponibilidad, caudal ecológico, obras y uso de agua.
\item 600 kW hidráulicos supera umbral de 500 kW para ruta de concesión.
\end{itemize}
\medskip\small Eventual SEIN 2046--2076: nuevos estudios y aprobación.\par
\alert{No hay permisos ni conexión obtenidos.}
\end{frame}
\begin{frame}{Costos y decisión recomendada}
\begin{columns}\begin{column}{.54\textwidth}\includegraphics[width=\linewidth,height=4.8cm,keepaspectratio]{figuras/costos-r03.png}\end{column}
\begin{column}{.43\textwidth}\small
Inversión estimada: \textbf{S/\dato{CapexM} millones}.\par
Precios de partidas estimados; sensibilidad \(\pm35\%\).\par\medskip
Escenarios financieros desfavorables.\par\medskip
Prioridades: aforos, topografía, padrón, placas y reutilización.\par
\alert{Comparar desarrollo por etapas antes de invertir.}
\end{column}\end{columns}\end{frame}
\begin{frame}{Conclusión: diseño académico con decisiones sustentadas}
\begin{itemize}
\item La demanda y el recurso se verifican antes de escoger capacidad.
\item La potencia nominal no está disponible permanentemente.
\item El diseño integra pérdidas, regulación, respaldo y seguridad.
\item Las fuentes, hipótesis y limitaciones quedan trazables.
\end{itemize}
\medskip\large Próximo paso: \textbf{medir y validar las bases del diseño}.\par
\small Memoria completa, ocho láminas y cálculos reproducibles R03.
\end{frame}
\end{document}
